\section*{Capítulo 10}

\solucion{10.1} Activos: $41.1+482.3+467.6=991.0$. Pasivos: $3{,}824.5+517.2+607.8+740.6+4.9=5{,}695.0$. Neto: $991.0-5{,}695.0=-4{,}704.0$ millones de soles. En dólares: $4{,}704.0/3.367=1{,}397.1$ millones (pasivo).

\solucion{10.2} $3{,}765.6/3.361-3{,}765.6/3.765=1{,}120.4-1{,}000.2=120.2$ millones. La Nota 14(c) reporta 119.3. La diferencia se debe a que la nota usa los tipos de cambio contables de cierre y a los movimientos del saldo durante el año, mientras que el cálculo usa los promedios del BCRP sobre un saldo constante.

\solucion{10.3} $2.901\times(80.00-92.60)=2.901\times(-12.60)=-36.6$ millones de dólares.

\solucion{10.4} Porque la deuda es a tasa fija y no hay riesgo de flujo de caja; el riesgo es de valor razonable y de refinanciación. Además, el componente dominante del costo es el spread de crédito, que ningún derivado de tasas puede fijar. La cobertura posible (tasa base de la refinanciación del CESCE) es parcial y debe esperar a que se conozcan monto y fecha.

\solucion{10.5} No es una variable de mercado con un precio que suba o baje: es una escasez de caja frente a obligaciones. Pero los derivados pueden exigir caja (primas, márgenes, colateral). Una empresa sin liquidez debe elegir instrumentos sin pago inicial y con colateral limitado; de lo contrario, la cobertura podría causar la crisis de liquidez que busca evitar.

\solucion{10.6} Suma impactos de escenarios distintos y de naturaleza distinta: el del tipo de cambio es un resultado contable, el del crudo un cambio de valor de inventario y el de tasas un cambio de valor razonable de bonos que no afecta la caja. Además, supone que todos los escenarios adversos ocurren a la vez, sin considerar correlaciones; y el de tasas se mide con signo arbitrario (una subida de tasas \emph{reduce} el valor de mercado de la deuda). Agregar riesgos exige un modelo conjunto.

\solucion{10.7} Protege el valor del inventario de crudo: el swap y el collar ganan cuando el WTI baja. Pero el margen de refinación se comprime 5 puntos y el inventario de productos (US\$ 313.5 millones) pierde cerca de 62.7 millones ($0.20\times313.5$), que ninguna cobertura sobre crudo compensa. Es el riesgo de \emph{crack} que el informe deja para la fase~2, con swaps de diésel y gasolina.
